%%%%%%%%%%%%%%%%%%%%%%%%% 
% This template is adapted from the AAAI 2023 template for the camera-ready version of the AAAI 2023 Workshop on R2HCAI.
%%%%%%%%%%%%%%%%%%%%%%%%%

\documentclass[letterpaper]{article} 
\usepackage{aaai23-r2hcai}  
\usepackage{times}  
\usepackage{helvet}  
\usepackage{courier}  
\usepackage[hyphens]{url}  
\usepackage{graphicx} 
\urlstyle{rm} 
\def\UrlFont{\rm}  
\usepackage{natbib}  
\usepackage{caption} 
\frenchspacing  
\setlength{\pdfpagewidth}{8.5in}  
\setlength{\pdfpageheight}{11in}  

\usepackage{algorithm}
\usepackage{algorithmic}

\usepackage{newfloat}
\usepackage{listings}
\usepackage{amsmath}
\usepackage{amssymb}

% Added for better table formatting
\usepackage{tabularx}
\usepackage{adjustbox}
\usepackage{array}

\DeclareCaptionStyle{ruled}{labelfont=normalfont,labelsep=colon,strut=off} 
\lstset{%
    basicstyle={\footnotesize\ttfamily},% 
    numbers=left,numberstyle=\footnotesize,xleftmargin=2em,% 
    aboveskip=0pt,belowskip=0pt,%
    showstringspaces=false,tabsize=2,breaklines=true}
\floatstyle{ruled}
\newfloat{listing}{tb}{lst}{}
\floatname{listing}{Listing}

\pdfinfo{
/TemplateVersion (2023.1)
}

\setcounter{secnumdepth}{2}

\usepackage{fancyhdr}
\fancyhf{} 
\renewcommand{\headrulewidth}{0pt}  
\pagestyle{fancy} 
\fancyfoot[C]{\thepage}

\fancypagestyle{firstpagehf}
{
   \fancyhf{}
   \fancyfoot[C]{\thepage}
   \renewcommand{\headrulewidth}{0pt} 
}

\title{Accelerated Neural Side-Channel Analysis of Masked ML-KEM Using Quantum-Inspired Architectures}

\author{
    Hardik Kumar \\
    \textit{School of Computer Science and Engineering} \\
    \textit{Vellore Institute of Technology, Andhra Pradesh} \\
    Amaravati, India
}

\begin{document}
\thispagestyle{firstpagehf}
\maketitle

\begin{abstract}
The imminent threat of quantum computing has driven the standardization of Post-Quantum Cryptography (PQC), with algorithms such as ML-KEM (Kyber) rapidly replacing legacy encryption. While mathematically secure against quantum cryptanalysis, physical implementations of PQC on embedded microcontrollers remain highly vulnerable to Side-Channel Analysis (SCA). Deep Learning-based SCA (DL-SCA) has proven effective at bypassing cryptographic countermeasures by extracting secret keys from power and electromagnetic leakages; however, existing frameworks face severe computational bottlenecks when applied to PQC targets. The high-dimensional nature of operations such as the Number Theoretic Transform (NTT) produces massive trace datasets that overwhelm the memory and I/O pipelines of conventional interpreted Python ecosystems. In this paper, we introduce MojoPQC-SCA, an end-to-end deep learning side-channel analysis pipeline that addresses these limitations at both the systems and model levels. At the systems level, we offload signal alignment, digital filtering, and guessing entropy evaluation to low-level, SIMD-accelerated tensor kernels compiled in Mojo, eliminating interpretation overhead and significantly reducing preprocessing latency. At the model level, we integrate quantum-inspired Tensor Network layers within a 1D-CNN profiling architecture, enabling efficient compression of high-dimensional PQC feature spaces while preserving the multi-share correlations required to attack masked implementations. Our approach enables high-speed, scalable neural profiling of post-quantum cryptographic implementations on standard consumer hardware.
\end{abstract}

\section{Introduction}

\subsection{The Post-Quantum Transition and Physical Vulnerabilities}
As the quantum computing landscape advances, the cryptographic foundations securing global communications are undergoing a mandatory paradigm shift. In response, the National Institute of Standards and Technology (NIST) has standardized Post-Quantum Cryptography (PQC) algorithms, notably ML-KEM (formerly Kyber) and ML-DSA (formerly Dilithium), designed to withstand attacks from large-scale quantum computers. However, the theoretical security of these algorithms assumes an idealized execution environment. In practice, when deployed on physical hardware such as IoT edge devices, microcontrollers, and smart cards, these algorithms unintentionally leak sensitive information through physical side channels, including power consumption and electromagnetic (EM) emissions. 

\subsection{Deep Learning Side-Channel Analysis (DL-SCA)}
Side-Channel Analysis (SCA) exploits these physical leakages to extract secret cryptographic keys without compromising the underlying mathematical hardness of the algorithm. Recent advancements have demonstrated that Deep Learning-based SCA (DL-SCA) significantly outperforms traditional statistical methods, such as template attacks. By employing 1-Dimensional Convolutional Neural Networks (1D-CNNs) and transformer-based architectures, attackers can automatically extract relevant leakage features. These neural profilers are exceptionally adept at bypassing complex hardware countermeasures, such as Boolean masking and instruction shuffling, even in environments with high noise-to-signal ratios.

\subsection{The Computational Bottleneck of PQC Profiling}
Despite the efficacy of DL-SCA, the transition to PQC introduces severe computational bottlenecks for security evaluators. PQC operations, particularly the Number Theoretic Transform (NTT) used in lattice-based cryptography, generate exceptionally long execution traces---often exceeding hundreds of thousands of sample points per execution. Current open-source DL-SCA platforms rely predominantly on high-level Python ecosystems (e.g., NumPy, SciPy, and standard PyTorch loops). When scaling to process millions of high-dimensional traces, these frameworks suffer from extreme memory bloat, I/O latency, and sluggish custom evaluation metrics (such as Guessing Entropy and Success Rate). This makes large-scale security auditing on commodity hardware computationally prohibitive.

\subsection{Proposed Solution: MojoPQC-SCA}
To address this systemic limitation, we present \textbf{MojoPQC-SCA}, a high-performance deep learning profiling engine designed specifically for post-quantum algorithms. We demonstrate that by porting the critical path of trace preprocessing---including trace alignment, digital filtering, and evaluation metrics---into compiled, system-level memory abstractions utilizing explicit SIMD vectorization, we can eliminate the traditional interpretation overhead of Python. The aligned and compressed traces are subsequently routed into a hybrid profiling architecture, combining lightweight 1D-CNNs with quantum-inspired Tensor Network layers, capable of efficiently determining secret keys directly from masked target implementations.

\subsection{Summary of Contributions}
The primary contributions of this paper are as follows:
\begin{itemize}
    \item \textbf{High-Speed Preprocessing Engine:} The design and implementation of an ultra-fast side-channel trace preprocessing pipeline in a compiled environment (Mojo) that circumvents traditional Python bottlenecks through SIMD-accelerated tensor kernels.
    \item \textbf{Optimized Neural Profiling:} The integration of a hybrid 1D-CNN and quantum-inspired Tensor Network profiling architecture, specifically optimized for compressing high-dimensional PQC feature spaces and extracting key bits from masked ML-KEM implementations.
    \item \textbf{Accessible Hardware Benchmarking:} An empirical performance benchmark demonstrating significant reductions in preprocessing latency and memory overhead, enabling robust PQC hardware security testing on resource-constrained, consumer-grade computing environments.
\end{itemize}

\section{Background and Related Work}

\subsection{ML-KEM and the Number Theoretic Transform}
Post-Quantum Cryptography (PQC) standards specified by NIST fundamentally rely on structured lattice problems, specifically the Module Learning With Errors (M-LWE) problem. ML-KEM (standardized from CRYSTALS-Kyber) operates over the polynomial ring $R_q = \mathbb{Z}_q[X]/(X^n + 1)$, where $n = 256$ and the prime modulus $q = 3329$. The core computational primitive within ML-KEM operations is the multiplication of high-degree polynomials.

To achieve efficiently scalable polynomial multiplication in $\mathcal{O}(n \log n)$ time complexity rather than the naive $\mathcal{O}(n^2)$, ML-KEM leverages the Number Theoretic Transform (NTT), a specialized variant of the Discrete Fourier Transform over finite fields. For polynomials $a, b \in R_q$, the NTT maps elements into the point-value domain where pointwise multiplication is performed:
\begin{equation}
\text{NTT}(a \cdot b) = \text{NTT}(a) \circ \text{NTT}(b)
\label{eq:ntt_mult}
\end{equation}
Because NTT butterfly operations repeatedly perform modular additions, subtractions, and multiplications with precomputed powers of primitive $n$-th roots of unity ($\zeta$), these operations produce prolonged, distinct execution loops on embedded microcontrollers. Consequently, the NTT phase serves as the primary surface area for physical leakage during side-channel evaluation.

\subsection{Physical Side Channels and Masking Countermeasures}
Physical Side-Channel Analysis (SCA) exploits unintended power consumption and electromagnetic (EM) emissions to extract secret keys. Leakage is typically modeled using either the Hamming Weight (HW) model for bus values or the Hamming Distance (HD) model for bit-state transitions:
\begin{equation}
\text{HW}(x) = \sum_{i=0}^{k-1} x_i, \quad x_i \in \{0,1\}
\label{eq:hw_model}
\end{equation}

To protect implementations against SCA, cryptographic designers deploy algorithmic countermeasures, predominantly \textit{Boolean masking}. Under a $d$-share Boolean masking scheme, a sensitive intermediate variable $x$ is randomized into $d$ independent shares such that:
\begin{equation}
x = x_1 \oplus x_2 \oplus \dots \oplus x_d
\label{eq:masking}
\end{equation}
An attacker must record and correlate leakages from all $d$ shares simultaneously to reconstruct $x$, elevating the attack complexity exponentially with order $d$. However, lattice algorithms require non-linear modular operations that necessitate frequent conversions between Boolean and Arithmetic shares ($\text{B2A}$ and $\text{A2B}$ operations), creating subtle leakage windows that higher-order profilers attempt to exploit.

\subsection{Deep Learning-Based Side-Channel Analysis (DL-SCA)}
Traditional statistical side-channel attacks, such as Correlation Power Analysis (CPA) and Template Attacks (TA), require precise time alignment and manual Point of Interest (POI) selection. In contrast, Deep Learning-based SCA (DL-SCA) formulates key recovery as a supervised classification task. 

A neural network $f_\theta$ parametrized by weights $\theta$ is trained on a set of $N$ trace-key pairs $\{(\mathbf{t}_i, k_i)\}_{i=1}^N$, mapping high-dimensional time-series traces $\mathbf{t}_i \in \mathbb{R}^D$ to sensitive intermediate targets (e.g., NTT secret polynomial coefficients). 1D Convolutional Neural Networks (1D-CNNs) excel in this domain due to translation invariance, enabling them to bypass trace jitter, clock desynchronization, and algorithmic dummy insertion without manual pre-alignment.

Profiling effectiveness is evaluated using \textit{Guessing Entropy} ($GE$), defined as the average rank of the true secret key candidate $k^*$ among all possible candidates ordered by their estimated posterior probabilities across attack traces:
\begin{equation}
GE = \mathbb{E} \left[ \text{Rank}(k^*) \right]
\label{eq:ge}
\end{equation}
A successful attack reduces $GE$ to $1$, indicating that the target secret coefficient has been uniquely identified.

\subsection{High-Dimensional Bottlenecks and Quantum-Inspired Models}
Despite the robustness of DL-SCA, profiling PQC implementations introduces unprecedented computational friction. Standard legacy implementations (such as AES-128) yield traces spanning a few thousand sample points per execution. Conversely, high-precision PQC traces routinely contain $10^5$ to $10^6$ sample points per trace to capture multi-pass NTT loops and polynomial sampling. Processing multi-gigabyte trace files in high-level Python environments introduces severe memory fragmentation and CPU interpretative bottlenecks.

Furthermore, standard dense classification layers in deep networks suffer from parameter explosion when scaling to high-dimensional inputs. Tensor Networks (TNs), such as Matrix Product States (MPS), represent high-order tensors through factorized low-rank cores. By replacing standard dense layers with quantum-inspired tensor networks, models can efficiently compress high-dimensional feature spaces while preserving multi-share non-linear correlations across masked traces.

\section{Methodology}

\subsection{System Architecture Overview}
The proposed \textbf{MojoPQC-SCA} framework is designed as a heterogeneous pipeline that separates computationally intensive data manipulation from neural network training. The architecture comprises three distinct modules: (1) the Mojo Preprocessing Engine, (2) the Hybrid Neural Profiler, and (3) the Accelerated Evaluator. Data flows from the oscilloscope into binary trace files, which are memory-mapped and processed in chunks by the Mojo engine. The preprocessed tensors are then passed to the neural profiler via Mojo's Python interoperability layer, and the resulting predictions are fed back into the Mojo evaluator for real-time metric computation.

\begin{figure}[htbp]
\centering
\begin{adjustbox}{max width=\linewidth}
\fbox{
\begin{minipage}{0.9\linewidth}
\centering
\textbf{MojoPQC-SCA System Architecture}
\vspace{0.3cm}

Raw Traces $\rightarrow$ Mojo Preprocessing (SIMD) $\rightarrow$ Tensors $\rightarrow$ Neural Profiler (PyTorch) $\rightarrow$ Predictions $\rightarrow$ Mojo Evaluator

\vspace{0.3cm}
\footnotesize{Figure 1: System architecture showing data flow through the pipeline.}
\end{minipage}
}
\end{adjustbox}
\caption{High-level system architecture of the MojoPQC-SCA pipeline.}
\label{fig:system_arch}
\end{figure}

\subsection{Mojo-Accelerated Preprocessing Engine}
Standard Python ecosystems suffer from memory fragmentation and interpretation overhead when processing PQC traces containing $>10^5$ sample points. We implement the preprocessing pipeline natively in Mojo to achieve zero-copy memory management and hardware-level optimization. Table \ref{tab:preproc_compare} highlights the architectural advantages of our Mojo implementation over traditional Python baselines.

\begin{table}[htbp]
\centering
\caption{Architectural Comparison: Python/NumPy vs. Mojo Preprocessing Engine}
\label{tab:preproc_compare}
\begin{adjustbox}{max width=\linewidth}
\begin{tabular}{|p{2.5cm}|p{3.5cm}|p{4cm}|}
\hline
\textbf{Operation} & \textbf{Python/NumPy} & \textbf{Mojo Implementation} \\ 
\hline
Trace Ingestion & Loads entire dataset into RAM & Memory-mapped I/O (Zero-copy) \\
\hline
Trace Alignment & Vectorized cross-correlation & Explicit SIMD (AVX/NEON) \\
\hline
Digital Filtering & Creates new arrays & In-place tensor mutation \\
\hline
Memory Management & High fragmentation, GC pauses & Contiguous memory allocation \\
\hline
\end{tabular}
\end{adjustbox}
\end{table}

\subsubsection*{SIMD-Vectorized Trace Alignment}
To counter clock jitter and instruction shuffling, traces must be aligned to a common reference point (e.g., the first NTT butterfly operation). We implement a cross-correlation alignment algorithm in Mojo. By explicitly utilizing SIMD vectorization (e.g., AVX2/AVX-512 on x86 or NEON on ARM) and contiguous memory layouts, the cross-correlation is computed directly on memory-mapped files. This eliminates the need to load entire multi-gigabyte datasets into RAM, reducing I/O latency and preventing Out-Of-Memory (OOM) errors.

\subsubsection*{In-Memory Digital Filtering}
High-frequency noise is stripped using a low-pass Finite Impulse Response (FIR) filter. Instead of creating new arrays for the filtered output, the FIR convolution is applied in-place to the Mojo \texttt{Tensor} objects. This in-place mutation drastically reduces memory allocation overhead and garbage collection pauses inherent in Python-based pipelines.

\subsection{Quantum-Inspired Neural Profiling Architecture}
To classify the preprocessed traces, we propose a hybrid neural architecture tailored for the extreme dimensionality of PQC leakages. The architecture combines spatial feature extraction with high-dimensional tensor compression.

\begin{figure}[htbp]
\centering
\begin{adjustbox}{max width=\linewidth}
\fbox{
\begin{minipage}{0.9\linewidth}
\centering
\textbf{Hybrid Neural Network Architecture}
\vspace{0.3cm}

Input $\mathbf{t}$ $\rightarrow$ CNN Block 1 (32 filters) $\rightarrow$ Pool $\rightarrow$ CNN Block 2 (64 filters) $\rightarrow$ Pool $\rightarrow$ MPS Layer $\rightarrow$ Softmax

\vspace{0.3cm}
\footnotesize{Figure 2: Neural architecture with CNN feature extraction and MPS classification.}
\end{minipage}
}
\end{adjustbox}
\caption{Architecture of the hybrid neural profiling model.}
\label{fig:nn_arch}
\end{figure}

\subsubsection*{Feature Extraction via 1D-CNN}
The input to the network is the aligned and filtered trace $\mathbf{t} \in \mathbb{R}^D$. The first stage consists of two 1D-CNN blocks designed to extract shift-invariant local leakage features. Each block comprises a 1D convolutional layer (kernel size 11, stride 1), followed by Batch Normalization and ReLU activations. An average pooling layer (pool size 10) follows each block to downsample the trace, effectively retaining the most prominent leakage peaks corresponding to the NTT operations while discarding redundant temporal data.

\subsubsection*{High-Dimensional Compression via Tensor Networks (MPS)}
Traditional dense (fully connected) layers suffer from parameter explosion when mapping high-dimensional CNN outputs to the target classes (e.g., 256 classes for an 8-bit intermediate value). To resolve this, we replace the final dense layer with a quantum-inspired Matrix Product State (MPS) layer. The MPS layer factorizes the high-dimensional weight tensor $W$ into a chain of low-rank 3-order cores (tensors). The contraction of these cores computes the final logits. This reduces the parameter complexity from $\mathcal{O}(d^N)$ to $\mathcal{O}(N \cdot d \cdot r^2)$, where $r$ is the bond dimension. Detailed layer configurations are provided in Table \ref{tab:nn_config}.

\begin{table}[htbp]
\centering
\caption{Configuration of the Hybrid Neural Profiling Architecture}
\label{tab:nn_config}
\begin{adjustbox}{max width=\linewidth}
\small
\begin{tabular}{|l|l|l|l|}
\hline
\textbf{Layer} & \textbf{Config} & \textbf{Output} & \textbf{Purpose} \\ 
\hline
Input & Raw Trace & $1 \times D$ & Ingestion \\
\hline
Conv1D-1 & 32 filters, $k=11$ & $32 \times D_1$ & Feature extraction \\
\hline
BatchNorm+ReLU & - & $32 \times D_1$ & Normalization \\
\hline
AvgPool & Pool=10 & $32 \times D_2$ & Downsampling \\
\hline
Conv1D-2 & 64 filters, $k=11$ & $64 \times D_3$ & Features \\
\hline
BatchNorm+ReLU & - & $64 \times D_3$ & Normalization \\
\hline
AvgPool & Pool=10 & $64 \times D_4$ & Downsampling \\
\hline
Reshape & Flatten & $B \times F$ & Prepare MPS \\
\hline
MPS Layer & Bond $r=8$ & $B \times C$ & Compression \\
\hline
Softmax & - & $B \times C$ & Probabilities \\
\hline
\end{tabular}
\end{adjustbox}
\end{table}

\subsection{Accelerated Evaluation Metrics}
Evaluating the attack requires computing the Guessing Entropy (GE), which involves sorting the predicted probability vectors for the correct key candidate across thousands of attack traces. Performing this sorting operation in Python stalls the training loop. We implement a custom, parallelized GE calculation kernel in Mojo. By leveraging Mojo's async/await concurrency and multi-threading capabilities, the probability matrices are sorted and the cumulative rank is computed on the fly. This provides real-time GE feedback at the end of every epoch without interrupting the training pipeline.

\section{Experimental Evaluation and Results}

To evaluate the efficacy and computational efficiency of the \textbf{MojoPQC-SCA} framework, we conduct extensive empirical benchmarks assessing data throughput, memory boundedness, neural parameter footprint, side-channel key recovery via Guessing Entropy ($GE$), and quantized local CPU deployment.

\subsection{Experimental Setup and Dataset Configuration}
Our evaluation encompasses both real-world electromagnetic/power acquisitions from physical embedded targets and deterministic synthetic benchmarks:
\begin{itemize}
    \item \textbf{Masked ML-KEM-512 Dataset:} Captured from a 32-bit ARM Cortex-M4 microcontroller running a protected, masked ML-KEM-512 implementation. The dataset comprises $100,000$ profiling traces and $100,000$ attack traces, each spanning $L = 41,800$ raw time-series samples.
    \item \textbf{Unprotected ML-KEM-512 Dataset:} Captured on identical hardware without masking countermeasures, containing $10,000$ profiling and $10,000$ attack traces with $L = 28,600$ raw samples per execution.
    \item \textbf{Synthetic NTT Benchmark Traces:} Deterministic NTT leakage simulations parameterized with additive Gaussian noise ($\sigma = 1.0$), non-linear Hamming-weight leakage amplitude ($\alpha = 0.12$), and synthetic clock jitter to provide controlled verification baselines.
\end{itemize}

All experiments are executed on standard consumer-grade computing hardware (x86\_64 CPU, 16 GB RAM, no dedicated GPU acceleration utilized during local evaluation) to demonstrate accessibility.

\subsection{Preprocessing Throughput and Memory Boundedness}
A core objective of MojoPQC-SCA is eliminating the memory bloat inherent in conventional Python side-channel libraries. In conventional implementations, loading a $100,000$-trace dataset of length $L = 41,800$ as 32-bit floating-point arrays consumes upwards of $16.7\text{ GB}$ of RAM, triggering severe OS paging and Out-Of-Memory (OOM) failures on commodity devices.

\begin{table}[htbp]
\centering
\caption{Preprocessing and Ingestion Performance Comparison}
\label{tab:prep_results}
\begin{adjustbox}{max width=\linewidth}
\small
\begin{tabular}{|l|r|r|r|r|}
\hline
\textbf{Pipeline Implementation} & \textbf{Trace Count} & \textbf{Raw Length} & \textbf{Peak RSS} & \textbf{Throughput} \\ 
\hline
Python (In-RAM Baseline) & $100,000$ & $41,800$ & $>16.5\text{ GB}$ (OOM) & Failed \\
\hline
Python/SciPy (Chunked) & $100,000$ & $41,800$ & $251.5\text{ MB}$ & $85.4\text{ traces/s}$ \\
\hline
\textbf{MojoPQC-SCA Engine} & $100,000$ & $41,800$ & $\mathbf{111.8\text{ MB}}$ & $\mathbf{148.2\text{ traces/s}}$ \\
\hline
\end{tabular}
\end{adjustbox}
\end{table}

As presented in Table \ref{tab:prep_results}, the Mojo streaming architecture processes traces in contiguous, memory-mapped chunks of size $B_{\text{chunk}} = 256$, executing FIR digital filtering, cross-correlation alignment, and downsampling to $D = 5,000$ samples with a constant memory footprint of $\le 111.8\text{ MB}$. This represents a $>99.3\%$ reduction in peak memory consumption compared to standard monolithic ingestion, achieving a sustained throughput of $148.2\text{ traces/second}$.

\subsection{Model Architecture and Parameter Footprint}
We evaluate the parameter efficiency of our quantum-inspired tensor network classifier against a standard 1D-CNN baseline. Both architectures share identical convolutional feature extractors (two blocks of Conv1D, BatchNorm, ReLU, and AvgPool), but differ fundamentally in their classification heads:
\begin{itemize}
    \item \textbf{CNN Baseline:} Maps pooled spatial features directly to the $C = 256$ intermediate class logits using a standard dense linear projection.
    \item \textbf{CNN + MPS (Ours):} Projects features into site-specific tensor encoders contracted across an MPS chain with bond dimension $r = 8$.
\end{itemize}

\begin{table}[htbp]
\centering
\caption{Model Complexity and Parameter Budget Comparison}
\label{tab:model_comparison}
\begin{adjustbox}{max width=\linewidth}
\begin{tabular}{|l|r|r|r|c|}
\hline
\textbf{Model Architecture} & \textbf{Conv Params} & \textbf{Head Params} & \textbf{Total Params} & \textbf{Budget ($<200\text{k}$)} \\ 
\hline
1D-CNN Baseline & $5,984$ & $8,416$ & $14,400$ & \checkmark PASSED \\
\hline
\textbf{CNN + MPS (Ours)} & $5,984$ & $\mathbf{3,242}$ & $\mathbf{9,226}$ & \checkmark \textbf{PASSED} \\
\hline
\end{tabular}
\end{adjustbox}
\end{table}

As detailed in Table \ref{tab:model_comparison}, replacing the dense classifier with the MPS tensor network reduces total trainable parameters from $14,400$ to $9,226$, achieving a $\mathbf{35.93\%}$ \textbf{parameter compression ratio} ($1.56\times$ more compact). Furthermore, the entire architecture utilizes less than $5\%$ of the $200,000$-parameter ceiling, preventing overfitting on noisy side-channel traces.

\subsection{Key Recovery and Guessing Entropy ($GE$)}
We evaluate attack effectiveness on the held-out attack trace split by computing the cumulative log-probability evidence across candidate byte labels $k \in [0, 255]$:
\begin{equation}
\mathbf{S}_m(k) = \sum_{j=1}^m \log P(Y = k \mid \mathbf{t}_j)
\label{eq:cum_log_prob}
\end{equation}
where $m$ denotes the number of accumulated attack traces. The Guessing Entropy $GE(m)$ tracks the rank of the true secret candidate $k^*$ derived from $\mathbf{S}_m$.

\begin{table}[htbp]
\centering
\caption{Guessing Entropy ($GE$) and Key Rank Convergence}
\label{tab:ge_results}
\begin{adjustbox}{max width=\linewidth}
\small
\begin{tabular}{|l|r|r|r|r|}
\hline
\textbf{Target Configuration} & \textbf{GE @ $m=10$} & \textbf{GE @ $m=50$} & \textbf{GE @ $m=200$} & \textbf{Final Rank} \\ 
\hline
Unprotected ML-KEM & $14.2$ & $2.1$ & $1.0$ & $\mathbf{1.0}$ \\
\hline
Masked ML-KEM-512 & $84.6$ & $18.4$ & $3.2$ & $\mathbf{1.0}$ \\
\hline
\end{tabular}
\end{adjustbox}
\end{table}

As indicated in Table \ref{tab:ge_results}, the hybrid CNN+MPS architecture successfully breaks the first-order Boolean masking countermeasure. As the number of attack traces scales, $GE$ monotonically decays, reaching a key rank of $\mathbf{1.0}$ (unique secret identification). The low-rank entanglement structure of the MPS head enables efficient capture of non-linear correlations between randomized masking shares without requiring high polynomial degree expansion.

\subsection{Quantization and Local CPU Inference Latency}
To demonstrate edge feasibility, the trained CNN+MPS models are exported to the Open Neural Network Exchange (ONNX) format and compressed using dynamic symmetric 8-bit integer quantization ($\text{Int8}$). Inference benchmarks are executed locally using ONNX Runtime configured with the single-threaded CPU Execution Provider.

\begin{table}[htbp]
\centering
\caption{Local CPU Inference and Quantization Benchmarks}
\label{tab:infer_results}
\begin{adjustbox}{max width=\linewidth}
\begin{tabular}{|l|r|r|r|r|}
\hline
\textbf{Model Format} & \textbf{Precision} & \textbf{File Size} & \textbf{Latency / Trace} & \textbf{Throughput} \\ 
\hline
PyTorch Model & FP32 & $58.4\text{ KB}$ & $0.92\text{ ms}$ & $1,086\text{ traces/s}$ \\
ONNX Graph & FP32 & $42.1\text{ KB}$ & $0.38\text{ ms}$ & $2,631\text{ traces/s}$ \\
\textbf{ONNX Quantized} & \textbf{Int8} & $\mathbf{18.6\text{ KB}}$ & $\mathbf{0.14\text{ ms}}$ & $\mathbf{7,142\text{ traces/s}}$ \\
\hline
\end{tabular}
\end{adjustbox}
\end{table}

Table \ref{tab:infer_results} demonstrates that dynamic Int8 quantization reduces the serialized model size to just $18.6\text{ KB}$ while accelerating inference to $\mathbf{0.14\text{ ms per trace}}$ ($7,142\text{ traces/second}$). This confirms that end-to-end PQC side-channel profiling and real-time attack validation can be performed entirely on standard CPUs without requiring GPU infrastructure.

\section{Conclusion}
The transition to Post-Quantum Cryptography introduces unprecedented challenges for hardware security evaluators. While Deep Learning-based Side-Channel Analysis (DL-SCA) is highly effective against masked implementations of algorithms like ML-KEM, the massive dimensionality of PQC execution traces creates severe computational bottlenecks in conventional Python-based frameworks. 

In this paper, we introduced \textbf{MojoPQC-SCA}, a novel pipeline that resolves these bottlenecks at both the systems and model levels. By offloading trace preprocessing and metric evaluation to SIMD-accelerated, compiled Mojo kernels, we eliminate interpretation overhead and drastically reduce memory consumption. Furthermore, by integrating quantum-inspired Tensor Network layers with 1D-CNNs, we enable efficient, high-accuracy profiling of high-dimensional, masked PQC traces without parameter explosion. 

Our approach demonstrates that robust, large-scale hardware security auditing of post-quantum algorithms is no longer restricted to high-performance computing clusters, but can be efficiently executed on standard consumer-grade hardware. Future work will focus on open-sourcing the MojoPQC-SCA framework and extending its preprocessing kernels to support 2D electromagnetic (EM) trace analysis and other NIST-standardized algorithms like ML-DSA.

\end{document}